\scanpage{28}
% transcription
\deleted{\unsure{applications}}

4. \deleted{\unsure{les}} $n$-champs \inserted{sur $X$} de $Y$-profondeur $\geq i$ \quad $(0\leq i\leq n+2)$, $Y$
étant un fermé de $X$.

5. Formules de chgt de base en top. ordinaire \deleted{\unsure{ou|à}}
\margin{$n$-champs ind-finis constructibles etc \ldots}
ou en topologie étale \ldots\ (Il y en a au moins trois \deleted{\unsure{versions}} \inserted{types}
principaux \ldots)

6. \deleted{\unsure{Exemples}} \inserted{Théorèmes} du type de Lefschetz \deleted{\unsure{en}} en topologie
étale (schémas) ou dans l'analytique-complexe.

7. $n$-champs localement constants.
Un morphisme de topos \inserted{$f:X\to Y$} est une \inserted{$n$-}équivalence d'homotopie
sss pour tout $n$ champ loc ct. \inserted{$A$} sur $Y$,
il induit une $n$-équivalence
\[
\Gamma(Y,A)\mathrel{\overset{(n)}{\simeq}}\Gamma(X,f^*(A))
\]
On peut donner une version pour les types
d'homotopie ``complétés'' \unsure{pour|par} une catégorie de groupes,
à la Artin-Mazur ?

\deleted{\unsure{Ce}} [ Cela suggère qu'un ``type d'homotopie'' de topos
est incarné par la connaissance, pour tt $n$,
de la $(n+1)$-catégorie des $n$-champs loc. constants
dessus, avec les foncteurs de passage des \deleted{\unsure{champs}} $n$-champs aux $(n-1)$-champs sous-jacents et \unsure{inversem.|inversement}
Dans le cadre d'Artin-Mazur, il faudrait des
restrictions de constructibilité convenables sur les
$n$-champs avec lesquels on travaille \ldots

8. Structures algébriques sur des $n$-catégories et
$n$-champs : Gr-structure, structure de Picard,
de Picard-stricte, torseurs etc. Stabilité par images
directes et inverses de $n$-champs.

9. La $(n+1)$-catégorie des $n$-champs de Picard \inserted{stricts} sur $X$
est $(n+1)$-équivalente à la sous-catégorie pleine des
% archive sheet number: 27
